\documentclass[../main.tex]{subfiles}

\begin{document}

\chapter{Soal Evaluasi (Ujian Tengah Semester dan Ujian Akhir Semester)}

\begin{subcpmk}
  \item Mengukur derajat kompetensi kognitif (CPL01, CPL02) serta porsi kelengkapan sintesis argumentatif Logika Komputasi berdasar porsi Sub-CPMK terukur dari keseluruhan semester pembelajaran arsitektur fungsional matematis.
\end{subcpmk}

\section{Panduan Evaluasi Orientasi \textit{Outcome-Based Education (OBE)} }
Pada fase ini, performa \textit{learning objectives} mahapeserta akan didominasi kemampuan konversi logika fungsional abstrak dan membuktikannya ke tataran spesifikasi rancangan algoritmik.

Batas capaian standar minimal kelulusan evaluatif *Outcome* kompetensi rasional di Universitas Bale Bandung adalah mahasiswa diposisikan kompeten dan sukses menyelesaikan sintesis penguraian fungsi jika melampaui standar parameter akurasi \textit{passing grade} nilai kumulatif minimum \textbf{70\% tingkat ketepat-gunaan logika rasional struktur algoritmik dan teknik argumentasinya}. Instrumen disusun menjangkau ranah kognitif C4 \textit{(Menganalisis/Analitik sintaksis)} hingga C6 \textit{(Mencipta/Merumuskan relasi Bukti)}.

\section{Bagian I : Set Instrumen Soal Ujian Tengah Semester (UTS)}
*Cakupan:* Evaluatif Sub-CPMK komputan rasional 1.1 -- 1.3 dan Parsial komputasi abstraksi dasar pemetaan logis 2.1 (Materi Tatap Muka Bab 2 hingga Bab 9).
\textit{Waktu Alokasi:} 90 Menit rasional hitung.
\textit{Metode:} Studi Kasus Analitik Sintaks logika Kertas \& Papan Abstrak \textit{Tracing}.

\textbf{Soal Komputasi Analisis:}
\begin{enumerate}
    \item Definisikan perbedaan mutlak tipe var kalimat Proposisional komputasi vs Fungsi Predikat \textit{Sentence} komputasional relasional algoritma beserta $1$ ilustrasi terapan sintaksnya di *programming case* fungsional \textit{real condition check}! \textbf{(15 Poin, Relevansi: CPMK-1)}
    \item Buatlah \textit{Truth Table} / Tabel Kebenaran faktual komputasional iteratif lengkap dan teliti dari parameter fungsi arsitektur *router* sistem di bawah ini. Serta konklusikan tipe klasifikasi validitas logika komputasional fungsionalnya (Apakah berkelas Tautologi, Kontradiksi logikal *System Crash*, atau Kontingensi mutlak fungsional).
    Formula Algoritma Kausalitas: $[\,p \land (p \rightarrow q)\,] \rightarrow q$
    \textbf{(25 Poin, Relevansi: CPMK-1)}
    \item Substitusikan dan \textit{trace evaluatif validitas} formal spesifikasinya secara argumentasi aljabar silogistik murni \textit{Inference Rule Base} dari urutan \textit{Logs file Server System Condition} ini, kemudian konklusikan hasil putusan algoritmik sistem \textit{security firewall}-nya:
    \textit{(Server Logging Premise Rules):} 
    1. Jika Jaringan *Gateway* terdeteksi merusak Integritas (\textit{J}), \textit{Back-up Router} logikal Aktif memblokir \textit{(R)}. 
    2. Jika Sisi *Back-up Router* mengklaim status aktif memblokir (\textit{R}), Kecepatan Kompresi Sistem akan direpresentasikan dengan penalti Minus/Turun Signifikan nilainya (\textit{K}). 
    3. Analisis dan \textit{Tracing Test}: Bila dilaporkan dari monitor \textit{Admin System} validitas variabel Kecepatan \textit{TIDAK* TURUN* komputasi rasionalnya/logikal $K$ statusnya dinotasikan term $*$False* / $\sim K$}, apa status logis yang absolut wajib diputuskan oleh sistem menyangkut kondisi \textit{Integrity} ranah jaringan di awal kausal (\textit{J})? Terangkan hukum penarikan komputasinya yang relevan diterapkan sistem!
    \textbf{(35 Poin, Relevansi: CPMK-1 Algoritma Deductive Engine)}
    \item Rumuskan Spesifikasi Fungsional Relasional Database predikat dari aturan deskriptif algoritma alam ini: "Tidak satupun dari fungsi var sistem *Thread Komputatif CPU ($x$)* di *Node Alpha* boleh dikalkulasi parameter status *Overheating ($O$)* kecuali sistem pendeteksi \textit{Sensor Cooling ($S$)} juga mem-berstatus predikat *Failure ($F$)*!" Gunakan fungsi parameter simbol Kuantifikasional Eksponensial $\exists$ atau Universal $\forall$.
    \textbf{(25 Poin, Relevansi: CPMK-2)}
\end{enumerate}

\section{Bagian II : Set Instrumen Ujian Akhir Semester (UAS)}
*Cakupan:* Evaluasi komprehensif mengintegrasikan parameter Sub-CPMK 2.2 -- 4.3 (Logika Predikat Bersarang, Boolean, dan Pembuktian). Integrasi CPMK-2, 3, 4 Mutlak.
\textit{Metode:} Desain Sirkuit Komputasi dan Verifikasi Teorema Algoritma.

\textbf{Soal Komputasi Analisis Kombinatorial \& Verifikatif:}
\begin{enumerate}
    \item Optimasi *Sum of Products*. Diberikan Spesifikasi instruktur *Logic output terminal Truth Table System* $F(A,B,C,D)$ bernilai *True/Menyala/1* manakala kombinasi posisi sel *Minterms*-nya menyentuh urutan \textit{Index} $\Sigma (0, 2, 4, 6, 8, 10, 12, 14)$.
    Gunakan peta Karnaugh \textit{K-Map} resolusi beraturan *4-Variable Matrix* untuk men-\textit{simplify} reduksi *cost gate circuit*-nya seoptimal ringkas rasional mungkin. \textbf{(20 Poin, Relevansi Sub-CPMK 3.2)}
    \item Konversi spesifikasi *Optimized Form* dari hasil parameter fungsi akhir reduksian nomor 1, kemudian desain gambaran arsitektur perkabelan *Gates Circuit*/Skemanya memakai modul Gerbang Utama dan Modul NOT!  \textbf{(20 Poin, Relevansi Sub-CPMK 3.3)}
    \item Uji komparator Metode *Proof* / \textit{Contrapositive Logical Rule}. Buktikan secara sah deduksi argumentatif sistematis \textit{Step Trace Verification}: "JIKA parametrik $5n + 3$ diarsiteki oleh kompilator sebagai var data nilai tipe Genap, MAKA argumen sistem menggaransi secara deduktif variabel numerik \textit{n}-nya wajib divaluasi setara bilangan berlogika ganjil!". \textbf{(30 Poin, Relevansi Sub-CPMK 4.1)}
    \item Gunakan Induksi Komputasi iteratif struktural (\textit{Mathematical Induction Algorithm Loop Rules}), konfirmasikan kepada *Reviewer System Programmer* bahwa formulasi struktur komputatif kalkulasi rasional $1^3 + 2^3 + 3^3 + \dots + n^3 = \left[\frac{n(n+1)}{2}\right]^2$ selalu bernilai absolut Benar bagi tiap *assignation assignment loop iterator index integer* rasional mutlak $(n \geq 1)$. Uraikan fase Basis \& Validasi Langkah transisinya secermat dan teratur!. \textbf{(30 Poin, Relevansi Sub-CPMK 4.2)}
\end{enumerate}

\begin{checklist}
  \item Berhasil melalui standar rubrik penilaian minimum pemahaman logikal akurasi $70$.
  \item Mampu mensintesis pengaplikasian komputasi spesifik \textit{System Analysis Boolean logic} untuk fondasi memprogram maupun mendesain elektronika.
\end{checklist}

\end{document}
